% =====================================================================
% ctex 中文论文最小可编译骨架
% 编译: latexmk -xelatex ctex_chinese.tex     (必须 XeLaTeX 或 LuaLaTeX!)
%       pdfLaTeX 不能直接排中文, 会字体报错
% 国标参考文献: 用 gbt7714 包, 条目须带 langid 区分中英文
% =====================================================================
\documentclass[UTF8,a4paper,12pt]{ctexart}
% 备选: \documentclass[UTF8]{article} + \usepackage{ctex}

\usepackage{amsmath,amssymb}
\usepackage{graphicx}
\usepackage{booktabs}     % 三线表
\usepackage[hidelinks]{hyperref}

% 国标文献 (需要 gbt7714 包与 .bib):
% \usepackage[backend=biber,style=gb7714-2015]{biblatex}  % biblatex 路线
% 或 bibtex 路线: \usepackage{gbt7714} + \bibliographystyle{gbt7714-numerical}

\title{中文论文标题}
\author{第一作者\quad 第二作者}
\date{\today}

\begin{document}

\maketitle

\begin{abstract}
这是中文摘要。用一段话讲清研究问题、方法、主要结果与贡献。

\noindent\textbf{关键词：} 关键词一；关键词二；关键词三
\end{abstract}

\section{引言}
正文开始。引用示例~\cite{ref1}。见图~\ref{fig:arch} 与式~\eqref{eq:main}。

\section{方法}
\begin{equation}\label{eq:main}
  y = f(x) + \epsilon
\end{equation}

\begin{table}[t]
  \centering
  \caption{实验结果对比}
  \label{tab:res}
  \begin{tabular}{lcc}
    \toprule
    方法 & 指标A & 指标B \\
    \midrule
    基线 & 0.80 & 0.75 \\
    本文 & 0.88 & 0.83 \\
    \bottomrule
  \end{tabular}
\end{table}

\begin{figure}[t]
  \centering
  % \includegraphics[width=0.6\textwidth]{fig1.pdf}
  \rule{0.6\textwidth}{3cm}   % 占位框, 有图时替换
  \caption{系统架构图}
  \label{fig:arch}
\end{figure}

\section{结论}
总结。

% --- 参考文献 -----------------------------------------------------
% 国标 bibtex 路线 (注释掉下面 thebibliography 后启用):
% \bibliographystyle{gbt7714-numerical}
% \bibliography{refs}
% 演示用内联:
\begin{thebibliography}{1}
\bibitem{ref1} 张三, 李四. 论文标题[J]. 期刊名, 2024, 1(1): 1-10.
\end{thebibliography}

\end{document}
